% 编译提示：使用 XeLaTeX 编译
% 需要 ctex 宏包支持中文

\documentclass[a4paper,9pt]{article}
\usepackage{amsmath}
\usepackage{amssymb}
\usepackage{geometry}
\usepackage{ctex}
\usepackage{xcolor}
\usepackage{multicol}
\usepackage{titlesec}
\usepackage{fancyhdr}
\usepackage{tcolorbox}
\usepackage{graphicx}
\usepackage[version=4]{mhchem}

\geometry{left=1cm,right=1cm,top=1cm,bottom=1.8cm}
\setlength{\columnsep}{0.8cm}
\setlength{\columnseprule}{0.4pt}
\renewcommand{\columnseprulecolor}{\color{gray!50}}

\definecolor{sectioncolor}{RGB}{0,102,204}
\definecolor{subsectioncolor}{RGB}{0,153,153}
\definecolor{keycolor}{RGB}{204,102,0}
\definecolor{derivcolor}{RGB}{100,149,237}
\definecolor{boxbg}{RGB}{240,248,255}
\definecolor{keybg}{RGB}{255,248,240}

\titleformat{\section}{\normalfont\large\bfseries\color{sectioncolor}}{\thesection}{0.5em}{}
\titlespacing*{\section}{0pt}{1ex plus 0.5ex minus 0.2ex}{0.8ex plus 0.2ex}
\titleformat{\subsection}{\normalfont\normalsize\bfseries\color{subsectioncolor}}{\thesubsection}{0.5em}{}
\titlespacing*{\subsection}{0pt}{0.8ex plus 0.3ex minus 0.1ex}{0.5ex plus 0.1ex}

\setcounter{secnumdepth}{4}
\setcounter{tocdepth}{4}

\newenvironment{keypoint}
{\par\vspace{0.3em}\noindent\begin{tcolorbox}[colback=keybg,colframe=keycolor,boxrule=0.5pt,arc=2pt,left=2pt,right=2pt,top=2pt,bottom=2pt,boxsep=0pt]\noindent\textcolor{keycolor}{\textbf{关键点：}}}
{\end{tcolorbox}\par\vspace{0.3em}}

\newenvironment{examplebox}
{\par\vspace{0.3em}\noindent\begin{tcolorbox}[colback=boxbg,colframe=derivcolor,boxrule=0.5pt,arc=2pt,left=2pt,right=2pt,top=2pt,bottom=2pt,boxsep=0pt]\scriptsize\noindent\textcolor{derivcolor}{\textbf{示例：}}\par\setlength{\parskip}{0.1ex}\setlength{\itemsep}{0pt}\setlength{\parsep}{0pt}}
{\end{tcolorbox}\par\vspace{0.3em}}

\pagestyle{fancy}
\fancyhf{}
\fancyfoot[C]{\scriptsize4-\thepage}
\renewcommand{\headrulewidth}{0pt}
\renewcommand{\footrulewidth}{0.4pt}
\renewcommand{\footrule}{\hbox to\headwidth{\color{gray!50}\leaders\hrule height \footrulewidth\hfill}}

\title{\vspace{-2em}\Large\bfseries\color{sectioncolor} Chapter 4 Collections - 振动、电子与同位素效应（Lect.5）\vspace{-1em}}
\author{}
\date{\today}

\begin{document}
\maketitle
\thispagestyle{fancy}
\begin{multicols}{2}
	\scriptsize{\tableofcontents}

	\section{振动配分函数}

	\subsection{简谐振子模型}
	振动自由度用简谐振子近似，能级为
	\[
		E_v=\left(v+\frac12\right)h\nu_0,
		\qquad v=0,1,2,\dots
	\]
	若相对基态取零点，则
	\[
		E_v=vh\nu_0.
	\]
	因此振动配分函数写成
	\[
		q'_{\mathrm{vib}}=\sum_{v=0}^{\infty} e^{-vh\nu_0/kT}
		=1+e^{-h\nu_0/kT}+e^{-2h\nu_0/kT}+\cdots.
	\]

	\begin{keypoint}
		振动与平动/高温转动不同：典型振动频率常在 $200\sim 4000\ \mathrm{cm^{-1}}$，室温下很多分子的振动能级间隔都大于 $kT\approx 200\ \mathrm{cm^{-1}}$，所以通常不能用积分近似，只能直接求和或利用级数闭式。
	\end{keypoint}

	\subsection{几何级数求和与 $\theta_{\mathrm{vib}}$}
	上式是公比为
	\[
		r=e^{-h\nu_0/kT}
	\]
	的几何级数，因此
	\[
		q'_{\mathrm{vib}}=\frac{1}{1-e^{-h\nu_0/kT}}.
	\]
	定义振动温度
	\[
		\theta_{\mathrm{vib}}=\frac{h\nu_0}{k}
	\]
	后，可写成
	\[
		q'_{\mathrm{vib}}=\frac{1}{1-e^{-\theta_{\mathrm{vib}}/T}}.
	\]

	\begin{keypoint}
		这里最好固定一种单位语言：若 $\nu_0$ 用 Hz，则
		\[
			\theta_{\mathrm{vib}}=\frac{h\nu_0}{k};
		\]
		若实验给的是波数 $\tilde\nu$（$\mathrm{cm^{-1}}$），则应写成
		\[
			\theta_{\mathrm{vib}}=\frac{hc\tilde\nu}{k},
		\]
		数值上常直接用
		\[
			\theta_{\mathrm{vib}}(\mathrm K)=\frac{\tilde\nu(\mathrm{cm^{-1}})}{0.6950}.
		\]
	\end{keypoint}

	若能量零点取在势阱底而非基态，则需乘上零点能修正：
	\[
		q_{\mathrm{vib}}=q'_{\mathrm{vib}}e^{-\varepsilon_0/kT},
		\qquad \varepsilon_0=\frac12 h\nu_0,
	\]
	故
	\[
		q_{\mathrm{vib}}=\frac{e^{-\theta_{\mathrm{vib}}/(2T)}}{1-e^{-\theta_{\mathrm{vib}}/T}}.
	\]

	\begin{examplebox}
		Ex.~9 的标准结果就是这两个版本：
		\[
			q'_{\mathrm{vib}}=\frac{1}{1-e^{-\theta_{\mathrm{vib}}/T}},
			\qquad
			q_{\mathrm{vib}}=\frac{e^{-\theta_{\mathrm{vib}}/(2T)}}{1-e^{-\theta_{\mathrm{vib}}/T}}.
		\]
		题目若说“以势阱底为零点”，就必须用后者；若说“以最低振动态为零点”，就用前者。
	\end{examplebox}

	\subsection{室温判断}
	当
	\[
		T\ll\theta_{\mathrm{vib}},
	\]
	只有基态显著占据，故
	\[
		q'_{\mathrm{vib}}\approx 1.
	\]
	例如 CO 的振动频率约 $2157\ \mathrm{cm^{-1}}$，所以
	\[
		\theta_{\mathrm{vib}}\approx \frac{2157}{0.6950}=3079\ \mathrm K,
	\]
	在 298 K 时几乎只有第一项起作用；但 I$_2$ 的振动频率较低，$\theta_{\mathrm{vib}}$ 与室温同量级，因此 $q'_{\mathrm{vib}}$ 会明显大于 1。

	\section{更复杂分子的振动}

	复杂分子的振动由\textbf{normal modes} 描述。若分子有 $N$ 个原子，则正常模数为
	\[
		3N-6\quad (\text{非线性}),
		\qquad
		3N-5\quad (\text{线性}).
	\]
	每个正常模都可看成独立简谐振子，第 $i$ 个模的配分函数为
	\[
		q_{\mathrm{vib},i}=\frac{e^{-\theta_{\mathrm{vib},i}/(2T)}}{1-e^{-\theta_{\mathrm{vib},i}/T}}.
	\]
	总振动配分函数为各模乘积：
	\[
		q_{\mathrm{vib}}=q_{\mathrm{vib},1}q_{\mathrm{vib},2}q_{\mathrm{vib},3}\cdots
	\]
	若某模有 $n$ 重简并，则同一因子出现 $n$ 次。

	\begin{examplebox}
		Ex.~11 的关键不是死套公式，而是先数正常模，再决定哪些模重复。比如 CO$_2$ 的弯曲模二重简并，因此总振动配分函数里弯曲模因子要平方。
	\end{examplebox}

	\section{电子配分函数}

	除氢原子外，电子能级一般没有简单量子力学模型可直接求出，通常要借助实验光谱数据。若电子能级有简并，则应写成
	\[
		q_{\mathrm{elec}}=\sum_j g_j e^{-\varepsilon_{\mathrm{elec},j}/kT}.
	\]
	若以电子基态为零点，
	\[
		q'_{\mathrm{elec}}=g_0+g_1 e^{-\varepsilon_{1}/kT}+g_2 e^{-\varepsilon_{2}/kT}+\cdots.
	\]
	大多数情况下，激发电子态远高于 $kT$，故
	\[
		q'_{\mathrm{elec}}\approx g_0.
	\]

	原子常用项符号 $^{2S+1}L_J$，该能级的简并度为
	\[
		g=2J+1.
	\]
	例如 Na 原子基态为 $^2S_{1/2}$，故 $g_0=2$；多数闭壳层分子电子基态不简并，因此常有 $q'_{\mathrm{elec}}=1$。O$_2$ 是例外，其基态为 $^3\Sigma_g^-$，故 $q'_{\mathrm{elec}}=3$。

	\begin{keypoint}
		电子配分函数通常数值不大，但电子基态能量本身在后面计算化学平衡常数时非常重要。因此要区分：
		\begin{itemize}
			\item 热激发部分常只看 $q'_{\mathrm{elec}}$；
			\item 绝对基态电子能量会进入反应的 $\Delta\varepsilon_0$。
		\end{itemize}
	\end{keypoint}

	\section{同位素取代}

	把分子中的某个原子换成其同位素，不改变核电荷数与电子数，因此通常可认为键长和力常数基本不变。于是转动与振动只通过约化质量 $\mu$ 改变：
	\[
		I=\mu R^2,
		\qquad
		B\propto \frac1I,
	\]
	所以
	\[
		\frac{B_A}{B_B}=\frac{\mu_B}{\mu_A}.
	\]
	振动频率满足
	\[
		\nu_0\propto \sqrt{\frac{k_{\mathrm{force}}}{\mu}},
	\]
	因此
	\[
		\frac{\nu_{0,A}}{\nu_{0,B}}=\sqrt{\frac{\mu_B}{\mu_A}}.
	\]

	\begin{keypoint}
		同位素比值题里，$B$ 可以是能量形式的转动常数，也可以是波数形式的转动常数；$\nu_0$ 也可以换成波数 $\tilde\nu$。但前后必须保持\textbf{同一种定义}，这样常数因子才会自动约掉。
	\end{keypoint}

	\begin{examplebox}
		Ex.~33/34 一类同位素交换平衡题，最核心的物理事实就是：同位素主要改动 $\mu$，进而改动 $\theta_{\mathrm{rot}}$ 和 $\theta_{\mathrm{vib}}$，而不会大幅改变势能曲线本身。
	\end{examplebox}

	\section{热力学函数可分离的起点}

	由于单分子配分函数可以分解为
	\[
		q=q_{\mathrm{trans}}q_{\mathrm{rot}}q_{\mathrm{vib}}q_{\mathrm{elec}},
	\]
	于是对不可分辨粒子体系
	\[
		A=-kT\ln\frac{q^N}{N!}
	\]
	可写成各部分和：
	\[
		A=A_{\mathrm{trans}}+A_{\mathrm{rot}}+A_{\mathrm{vib}}+A_{\mathrm{elec}}.
	\]
	通常把唯一的 $N!$ 因子归到平动部分，因此
	\[
		A_{\mathrm{trans}}=-NkT\ln q_{\mathrm{trans}}+kT(N\ln N-N),
	\]
	而其余部分都只是
	\[
		A_x=-NkT\ln q_x.
	\]

	同理，内能通式是
	\[
		U=NkT^2\left(\frac{\partial\ln q}{\partial T}\right)_V,
	\]
	因此后面每一部分都能单独计算。

	\begin{keypoint}
		从这一讲开始，思路就固定了：
		\[
			q_x \longrightarrow \ln q_x \longrightarrow U_x,S_x,C_{V,x}.
		\]
		后面 lect6--9 的大多数题目，本质上都是这条链在不同物理自由度上的重复。
	\end{keypoint}

	\section{最后速记}
	\begin{keypoint}
		Lect.~5 高频公式：
		\[
			q'_{\mathrm{vib}}=\frac{1}{1-e^{-\theta_{\mathrm{vib}}/T}},
			\qquad
			q_{\mathrm{vib}}=\frac{e^{-\theta_{\mathrm{vib}}/(2T)}}{1-e^{-\theta_{\mathrm{vib}}/T}},
		\]
		\[
			q_{\mathrm{vib}}=\prod_i q_{\mathrm{vib},i},
			\qquad
			q_{\mathrm{elec}}=\sum_j g_j e^{-\varepsilon_j/kT},
			\qquad
			q'_{\mathrm{elec}}\approx g_0,
		\]
		\[
			\frac{B_A}{B_B}=\frac{\mu_B}{\mu_A},
			\qquad
			\frac{\nu_{0,A}}{\nu_{0,B}}=\sqrt{\frac{\mu_B}{\mu_A}},
			\qquad
			q=q_{\mathrm{trans}}q_{\mathrm{rot}}q_{\mathrm{vib}}q_{\mathrm{elec}}.
		\]
	\end{keypoint}

\end{multicols}
\end{document}
